\begin{reviewcomment}{5}
	The comparison with related work should be strengthened by considering recent blockchain-based access-control architectures that combine smart contracts, SSI-oriented trust models, and cryptographic mechanisms for controlling access to sensitive claims.
	In particular, the paper should better position TPS-ZKPI with respect to the issues discussed in DOI: 10.1016/j.bcra.2024.100196, especially regarding verifier authorization, smart-contract accountability, off-chain trust assumptions, and practical access-control workflows.
\end{reviewcomment}

\begin{reviewresponse}{5}
	We added a paragraph to the Related Work section that positions TPS-ZKPI against Buccafurri et~al.~[\Loc{ref-buccafurri2024holder}] on the four points raised.
	\emph{Verifier authorization:} there the holder's claims are protected, whereas here the provider's policy is protected and any observer can verify the proof.
	\emph{Smart-contract accountability:} the contract records the context-bound outcome rather than metering the release of claims.
	\emph{Off-chain trust:} the TPS sees both the policy and the claims, which Table~3 makes explicit.
	\emph{Workflow:} holder-authorized presentation could be combined with TPS-ZKPI.
	We also cite recent zero-knowledge authentication work~[\Loc{ref-madine2025zero}] and a survey~[\Loc{ref-shashidhara2024promise}].

	\RevisionLocation{Section~3, fourth paragraph (\Loc{rw-ssi}); References.}
\end{reviewresponse}
